\hypertarget{commercial-latent-variable-methods-and-the-banking-question}{%
\section{Commercial latent-variable methods and the banking
question}\label{commercial-latent-variable-methods-and-the-banking-question}}

\label{exec13:latent-commercial}

A relationship manager can observe borrower risk that is absent from the
analyst's table. The same risk can influence a rate concession and the
subsequent credit loss. A commercial system that offers latent variables
is therefore a relevant alternative, but that phrase covers several
different services. An analyst may specify a measurement model and
estimate its parameters; a clustering procedure may construct unobserved
customer classes; a discovery procedure may retain causal relationships
compatible with hidden common causes; or a randomized experiment may
identify an effect for an unobserved compliance class. The proposed
product needs to choose among these services using the decision and
information available to the bank.

This comparison examines the declared GeNIe manual family 5.0.R2,
mutable BayesiaLab and Netica help, the Stata latent-class example
carrying a 2025 copyright and release-19 data reference, SAS/STAT 14.1
CALIS documentation, and Chapter 7 of the guide linked by Mplus as
Version 8. These are specific documentation observations. The public
copies were acquired after the 7 October 2026 literature cutoff; their
headers and links identify editions without establishing every page's
exact bytes at that cutoff. No commercial executable was run, and
documentation of one path cannot establish the presence or absence of
every other method in the vendor's ecosystem.
\citep{exec13_genie_readme,exec13_latent_genie_parameters,
exec13_latent_bayesia_data,exec13_latent_netica_structure,
exec13_latent_stata_lca,exec13_latent_sas_calis,exec13_latent_mplus}

\hypertarget{the-object-that-must-survive-a-discovery-interface}{%
\subsection{The object that must survive a discovery
interface}\label{the-object-that-must-survive-a-discovery-interface}}

For causal discovery with unmeasured common causes, a learned latent
class is different from an equivalence class of causal graphs. In the
no-selection ancestral-graph construction reviewed by
\citet[Section 3]{zhang2008}, a maximal ancestral graph represents
marginal independence and ancestral relations among observed variables.
A directed edge records ancestry in the underlying causal DAG; it need
not exclude additional hidden confounding. A bidirected edge records the
absence of ancestry in either direction and the presence of a latent
common-cause path. It does not identify a particular borrower-risk
factor, its dimension, or its measurement equation.

A partial ancestral graph, or PAG, represents a Markov equivalence class
of these graphs. Its invariant tails and arrowheads are shared across
the class, while a circle records an unresolved endpoint. Under the
stated Markov, faithfulness and graph assumptions,
conditional-independence information can constrain this object. A fitted
mixture component or a posterior probability of customer membership is a
different mathematical output. Neither can replace a circle endpoint
merely because its label contains the word ``latent.'' The cited
construction excludes selection; Zhang explicitly notes that selection
requires a broader ancestral-graph class and additional rules. A
loan-book analysis conditional on approval must retain that extra
question. \citep[pp.1442--1445]{zhang2008}

The proposed library consequently needs separate representations for a
supplied latent-variable model, a fitted measurement or mixture model, a
discovered graph class and an identified intervention query. Parameters,
class probabilities and graph endpoints have different meanings. The
separation also determines the correct baseline: a commercial mixture
fit should be compared with the same mixture specification, while a
proposed PAG operation should be compared with a discovery method under
the same latent and selection assumptions. The presence of a diagram is
insufficient to make these comparisons interchangeable.

\hypertarget{a-good-observational-fit-can-leave-the-lending-effect-unresolved}{%
\subsection{A good observational fit can leave the lending effect
unresolved}\label{a-good-observational-fit-can-leave-the-lending-effect-unresolved}}

Let \(A\) be a centered concession and \(Y\) a centered subsequent loss.
Consider two Gaussian structural models with independent disturbances.
In the first, \[
 A=\varepsilon_A,\qquad
 Y=\frac35 A+\varepsilon_Y,\qquad
 \operatorname{var}(\varepsilon_A)=1,\quad
 \operatorname{var}(\varepsilon_Y)=\frac{16}{25}.
\] In the second, \[
 A=U+\eta_A,\qquad Y=U+\eta_Y,\qquad
 \operatorname{var}(U)=\frac35,\quad
 \operatorname{var}(\eta_A)=\operatorname{var}(\eta_Y)=\frac25.
\] The unobserved \(U\) represents a stipulated risk component. Both
models have exactly the same observed covariance matrix, \[
 \Sigma_{(A,Y)}=
 \begin{pmatrix}1&3/5\\3/5&1\end{pmatrix},
\] and therefore the same centered bivariate Gaussian distribution.
Nevertheless, replacing the equation for \(A\) changes the mean of \(Y\)
by \(3/5\) per unit in the first model and by zero in the second. The
second model's observational regression slope remains \(3/5\); it
describes risk-related assignment, not an intervention response. These
two models establish nonidentification within any proposed model class
containing both. They do not enumerate every compatible model or
establish that the full identified set contains only two effects.
\citep{exec13_latent_examples}

For an explicitly illustrative monetary scale, let one concession unit
equal one percentage point and one centered-loss unit equal one thousand
currency units per borrower. A 25 basis point action then changes
expected loss by 150 currency units under the first model and zero under
the second. The equality of the observed likelihoods cannot choose
between those recommendations. Increasing the observational sample or
selecting the better optimizer cannot resolve an equality of
distributions. The calculation is exact project algebra with stipulated
units, not a fitted bank estimate or a vendor result.
\citep{exec13_latent_examples}

A risk score or inferred segment also needs a measurement argument
before it can close that gap. In a separate Gaussian example, take
independent unit-variance \(U,\varepsilon_A,\varepsilon_Y\), and let \[
 A=U+\varepsilon_A,\qquad Y=U+\varepsilon_Y,\qquad
 W=U+\varepsilon_W,
\] with independent proxy noise of variance \(s\). There is still no
causal effect of \(A\) on \(Y\). Given \(W\), the residual variance of
risk is \(v=s/(1+s)\). Consequently, \[
 \frac{\operatorname{cov}(A,Y\mid W)}
      {\operatorname{var}(A\mid W)}=
 \frac{v}{1+v}.
\] With \(s=1\), adjustment reduces the observational slope from \(1/2\)
to \(1/3\), but the causal effect remains zero. Only the stipulated
error-free limit removes this particular source of confounding. This
example does not rule out identification from appropriately designed
multiple proxies or a justified measurement model; it shows why the mere
availability of an inferred score is insufficient. If the score is
trained using post-offer outcomes, its information clock needs an
additional review. \citep{exec13_latent_examples}

\hypertarget{existing-engines-for-supplied-models-and-latent-classes}{%
\subsection{Existing engines for supplied models and latent
classes}\label{existing-engines-for-supplied-models-and-latent-classes}}

GeNIe's parameter-learning documentation starts with an existing network
whose structure is already defined. Its EM procedure can learn
parameters with missing observations and explicitly discusses random
initialization when latent variables are present. The source states that
parameter learning leaves structure fixed. It also makes a consequential
distinction: its continuous multivariate-Gaussian EM implementation
described on that page does not allow missing values. The proposal
cannot infer that every data-type/missingness combination is supported
from the general EM label. \citep{exec13_latent_genie_parameters}

The structure-learning introduction describes a separate interface over
selected data columns. In that declared manual, the routines accept
categorical variables; PC additionally accepts jointly Gaussian
continuous variables. Mixed categorical/continuous structure learning
requires discretization, and the manual states a missing-record
restriction with an exception for the Naive Bayes routine that
constructs its prescribed structure. These are documented restrictions
of the selected interface, not a claim about all current BayesFusion
products. An initial reuse route is therefore a supplied discrete latent
network with an explicit observation model and parameter-fitting
contract. Its fitted parameters must not be advertised as a discovery
result that has resolved hidden confounding.
\citep{exec13_latent_genie_structure,exec13_latent_genie_parameters}

BayesiaLab supplies substantive latent-representation alternatives. Its
Data Clustering documentation creates a new factor node whose states are
clusters. The ordinary construction uses a Naive Bayes structure with
the factor as parent of the manifest variables, random initialization,
and EM updates based on posterior class probabilities. Conditional
independence of the manifest variables given the class is part of that
representation. The documented Multinet alternative combines a Naive
structure with class-specific maximum-weight spanning trees and hard
class assignments during its alternating procedure. Thus the vendor
offers more than a single generic clustering option; the representation
and assignment policy belong in a matched comparison.
\citep{exec13_latent_bayesia_data}

Its heterogeneity option changes the objective toward segmentations that
increase a stated mutual-information index with the target. That can be
useful for describing different customer groups. An increase in that
index is neither an estimate of treatment-effect heterogeneity nor
evidence that the induced factor blocks every confounding path. The
adjacent Variable Clustering tool groups nodes using Arc Force, offers a
dendrogram, and can feed Multiple Clustering to construct synthetic
variables. It groups variables rather than immediately identifying a
customer's hidden risk. A proposed integration must retain which of
these operations was performed, its variables and target, soft versus
hard assignments, and the data used to construct the representation.
\citep{exec13_latent_bayesia_data,exec13_latent_bayesia_variables}

Netica's Learning Algorithms help describes counting, EM and conjugate
gradient methods for conditional probability tables. Counting uses cases
with definite values for the node and its parents; the page directs
latent or incomplete-data applications toward EM or gradient methods and
notes local optima. Its adjacent Structure Learning page describes TAN
learning around a selected diagnostic or predictive target, expressly
motivating outgoing target links by predictive dependence. This is a
useful modeling and prediction alternative. It does not supply a causal
interpretation of those target-to-predictor arrows. The structure page
refers to the version 5.00 introduction of TAN; that historical feature
description is not a verified inventory of a current Netica release.
\citep{exec13_latent_netica_algorithms,exec13_latent_netica_structure}

The Netica discussion writes an objective proportional to
\(p(D\mid N)p(N)\), where \(N\) denotes the candidate network and its
parameters. When the prior varies across candidates, this is a posterior
criterion, even though the accompanying discussion also uses
maximum-likelihood language. A matched adapter comparison must therefore
record the prior and existing probability tables as well as the reported
fit statistic. The product should compare the actual configured
objective; a shared method name does not establish numerical
equivalence. \citep{exec13_latent_netica_algorithms}

Stata's worked latent-class example is more explicit about the fitted
quantity. It supplies a two-class categorical latent variable through
\texttt{lclass(C 2)} and estimates class-specific intercepts for four
binary responses using \texttt{gsem}. The example reports marginal class
probabilities, posterior membership probabilities and hard assignments
obtained by thresholding those probabilities. These are distinct
summaries: the reported marginal proportions need not equal the
fractions of people assigned by the threshold rule. Neither summary is
automatically an observed pre-offer covariate or a randomized treatment
class. \citep[pp.2--6]{exec13_latent_stata_lca}

The same Stata example shows an EM initialization reaching its iteration
limit before full-model fitting converges. That is not evidence that the
final fit failed; it demonstrates why an integration should preserve
stage-specific diagnostics instead of interpreting the first convergence
message as the entire result. A prototype can reuse this supplied-model
workflow as a latent-class comparator while retaining the chosen class
count, starts, fitted probabilities, uncertainty and final convergence
status. No Stata execution is reported here.
\citep[pp.3--5]{exec13_latent_stata_lca}

SAS/STAT 14.1 CALIS provides another substantial alternative for
specified latent models. Its overview describes factor,
measurement-error, covariance, path and other structural-equation
specifications, several modeling languages, parameter constraints, fit
diagnostics and reusable output data sets. ML and GLS carry
multivariate-normal assumptions in the documented formulation. FIML
handles missing data under the stated missing-at-random assumption. The
source distinguishes Satorra--Bentler scaled inference from the separate
robust procedure that downweights outliers. These choices address
different statistical problems; selecting one does not validate the
causal assignment or make nonrandom observation ignorable. The new
product currently lacks this breadth of fitted-model workflow and should
justify any replacement rather than treating a new interface as
equivalent functionality. \citep{exec13_latent_sas_calis}

\hypertarget{mplus-is-also-a-causal-effect-comparator}{%
\subsection{Mplus is also a causal-effect
comparator}\label{mplus-is-also-a-causal-effect-comparator}}

The Mplus mixture chapter distinguishes measurement and structural
models, supports categorical latent classes and extensions with
continuous latent variables, and describes multiple starts for local
solutions. It documents options for missing data and complex sampling,
including strata, clusters and weights. These declarations identify
specific modeling alternatives; they do not imply that any chosen
default handles the bank's sampling or missingness law. A comparison
should reproduce the selected model and settings before measuring
resource use or decision quality.
\citep[pp.165--171]{exec13_latent_mplus}

More importantly for this proposal, Examples 7.23 and 7.24 describe
complier-average causal-effect estimation in randomized trials. In
Example 7.23, compliance class is observed in the treatment arm and
latent in the control arm. Training indicators constrain membership, the
noncomplier treatment slope is fixed at zero, and the complier slope is
estimated. The source identifies maximum likelihood with robust standard
errors as the default for this example. This is an existing documented
causal-effect use of a latent-class model. The proposal cannot claim
novelty merely from combining a latent variable with an intervention
target. \citep[pp.203--205]{exec13_latent_mplus}

A banking interpretation is a randomized assistance offer. In this
example, \(Y\) denotes net customer value in thousands of currency
units, so a positive effect is beneficial. Let \(Z\) denote the offer,
\(D\) receipt of assistance, and \(C\) a pre-existing compliance type.
For a declared one-sided design, suppose \(D=ZC\), \(Z\) is randomized
independently of potential outcomes and \(C\), and the offer has no
direct effect on the outcome apart from receipt. Consistency and a
positive first stage are also required. Under these restrictions, with
\(p_c=\Pr(C=1)\), the offer contrasts satisfy \[
 \Delta_D=p_c,\qquad
 \Delta_Y=p_c\,\operatorname{E}\{Y(1)-Y(0)\mid C=1\}.
\] The ratio identifies an effect for compliers, rather than the effect
of offering assistance to every eligible borrower. The experimental
design and the exclusion restriction supply the interpretation; a fitted
class alone does not.

An exact four-cell project law illustrates the distinction. Set
\(p_c=2/5\), assign the offer with probability \(1/2\), and define
\(Y=2+C+3D\) in thousands of currency units. The first stage is \(2/5\),
the offer effect is \(6/5\), and the receipt effect for compliers is
three. If an additional direct offer effect of \(1/5\) is present, the
offer effect becomes \(7/5\) and the same ratio becomes \(7/2\),
although the effect of receipt holding the offer fixed remains three.
Communication or reassurance accompanying an offer can therefore matter
to the model. The exact enumeration is a proposed-design diagnostic, not
an Mplus replication or an empirical estimate.
\citep{exec13_latent_examples}

Even a documented encoding deserves a numerical check. Example 7.24 sets
binary-indicator logit thresholds to \(+15\) and \(-15\) while
describing perfect class measurement. The smaller implied probability is
\(1/(1+e^{15})\), approximately \(3.0590\times10^{-7}\), rather than an
exact structural zero. This observation does not demonstrate an
empirical failure of Mplus. It means that a proposed adapter must
distinguish a deliberately near-deterministic measurement specification
from a hard membership constraint, including in comparisons with Example
7.23's training-variable route. \citep[pp.205--207]{exec13_latent_mplus,
exec13_latent_examples}

\hypertarget{what-to-reuse-and-what-the-prototype-must-add}{%
\subsection{What to reuse and what the prototype must
add}\label{what-to-reuse-and-what-the-prototype-must-add}}

\begin{longtable}{@{}P{.22\linewidth}P{.35\linewidth}P{.35\linewidth}@{}}
\caption{Commercial latent-variable paths and their proposed role.}
\label{exec13:latent-commercial-table}\\
\toprule
Inspected path & Documented object or operation & Proposed role and remaining check \\
\midrule
\endfirsthead
\toprule
Inspected path & Documented object or operation & Proposed role and remaining check \\
\midrule
\endhead
GeNIe parameter learning & EM parameters for an existing network; type-specific missingness limits & Reuse a supplied-model fit after checking the selected engine interface and observation model. \\
BayesiaLab clustering & Factor/class representation; Naive and Multinet alternatives & Retain as representation and segmentation baselines; require a separate causal measurement argument. \\
Netica learning & CPT estimation; separately documented predictive TAN structure & Compare matched CPT objectives and retain predictive graph semantics. \\
Stata gsem & Specified class count, class-specific responses and membership probabilities & Compare a supplied latent-class fit; preserve starts, stages and soft assignments. \\
SAS CALIS & Specified latent SEM, constraints, diagnostics and missing-data options & Reuse a model-fitting workflow where its assumptions match; preserve covariance and measurement restrictions. \\
Mplus mixture/CACE & Mixture parameters or a trial-specific complier effect & Include a causal-effect baseline; retain randomization, compliance restrictions and the complier target. \\
\bottomrule
\end{longtable}

The commercial systems provide capabilities the proposal has not
implemented: mature model specification, latent-class fitting,
diagnostic displays, specialized estimation choices and, in the Mplus
example, an experimental causal estimand. The proposed contribution is a
verifiable connection from the banking question to these operations,
with preserved assumptions, graph ambiguity, observation clocks and
monetary targets. That contribution remains a prototype obligation. A
generic CLI or MCP endpoint cannot create the missing measurement or
assignment evidence.

Direct adoption as a complete latent-confounded discovery service is
unsupported by this document set. That is a statement about the reviewed
evidence, not a finding that the vendors lack such services. Selected
commercial engines may instead be integrated for supplied-model or
experimental-effect operations. The BayesFusion and Bayesia deployment
distinctions reviewed earlier still apply; the selected edition, engine
interface and rights would accompany any actual trial. Porting a
proprietary implementation has no justification from public
documentation alone. \citep{exec13_smile_license,exec13_bayesia_api}

For the next inexpensive prototype comparison, use public or generated
cases with known information conditions. The two Gaussian worlds test
whether an interface invents a unique effect from indistinguishable
data; the proxy example tests whether a score is silently treated as
measured confounder truth; the randomized-offer law tests preservation
of offer, receipt and complier targets. A separate graph-class case
should test PAG endpoint roundtrips under an explicitly reviewed
no-selection profile. Selection requires its own case and graph
contract. These deterministic references can precede licensed vendor
execution and a prospectively registered sampling study. Any later
performance comparison must retain failed fits, initializations, budgets
and parameter constraints, and compare the same scientific operation
across engines.

Coverage also remains bounded. Edition-specific latent-discovery
capabilities of HUGIN, MathWorks, IBM Amos, LISREL and causaLens remain
unresolved in this comparison. They are candidates for further primary
technical review, with their information gaps retained in the
accompanying source register. No market-wide absence, cutoff-complete
census, installed interoperability, or paid-product superiority follows
from the present comparison. The immediate investment choice is to reuse
and test selected documented operations while keeping latent-confounded
discovery and effect identification explicit.
\label{exec13:latent-commercial-end}
